% GENERATED FILE. Edit canonical lessons, then run npm run generate:books.
% canonical-source-hash: fnv1a64:2dfa1bf566356ee2
% canonical-lessons: TA-C437-vanavilanku, TA-C437-caranalayam, TA-C437-itampeyar, TA-C437-vettaiyatu, TA-C437-inam

\chapter{Wildlife on the Move}
\label{ch:ta-wildlife-on-the-move}
\hlchaptermodality{437}

\begin{chapteropening}
\textbf{By the end of this chapter:} \emph{I can talk about wildlife and sanctuaries, species, and animals migrating and hunting.}

Complete the last lesson of chapter 437: I can talk about wildlife and sanctuaries, species, and animals migrating and hunting.
\end{chapteropening}

\section[\texorpdfstring{\ta{வனவிலங்கு} (vaṉavilaṅku)}{வனவிலங்கு (vaṉavilaṅku)}]{\ta{வனவிலங்கு} (vaṉavilaṅku) — wildlife; a wild animal}

\label{lesson:TA-C437-vanavilanku}



\begin{quote}
Before the new one: say the Tamil for a seashell, then the Tamil for coral.
\end{quote}

\subsection*{You'll want to know: \ta{வனவிலங்கு}}
\textbf{\ta{வனவிலங்கு}} — \emph{vaṉavilaṅku} — \textquotedblleft{}wildlife; a wild animal\textquotedblright{}.

\begin{quote}
\ta{இந்தக்} \ta{காட்டில்} \ta{நிறைய} \ta{வனவிலங்குகள்} \ta{இருக்கின்றன}.
\end{quote}

\begin{quote}
\emph{intak kāṭṭil niṟaiya vaṉavilaṅkukaḷ irukkiṉṟaṉa.}
\end{quote}

\begin{quote}
There are many wild animals in this forest.
\end{quote}

\ta{வனம்}, a forest, + \ta{விலங்கு}, an animal.

\begin{cousinweb}
\textbf{\ta{வனவிலங்கு}} (\emph{vaṉavilaṅku}) is built from \ta{வனம்}, \textquotedblleft{}a forest\textquotedblright{}, borrowed from Sanskrit \emph{vana} (CDIAL 11258), and native \ta{விலங்கு}, \textquotedblleft{}an animal\textquotedblright{}.
\end{cousinweb}

\begin{grammarlens}[title={What you've built}]
One more word for getting around.
\end{grammarlens}

\subsection*{Your turn}
\begin{itemize}
\raggedright
  \item \emph{Say it:} \emph{vaṉavilaṅku}
  \item \emph{Say it:} \emph{vaṉavilaṅku}, once more
  \item \emph{Say it:} \emph{intak kāṭṭil niṟaiya vaṉavilaṅkukaḷ irukkiṉṟaṉa.} — There are many wild animals in this forest.
  \item \emph{Recall:} say the Tamil for a seashell, then the Tamil for coral, then say \emph{vaṉavilaṅku} again
\end{itemize}

\begin{tcolorbox}[breakable,colback=teal!4,colframe=teal!35!black,title={Before you move on}]
What does \ta{வனவிலங்கு} mean? (\textquotedblleft{}Wildlife; a wild animal\textquotedblright{}.) Say it once more.
\end{tcolorbox}
\section[\texorpdfstring{\ta{சரணாலயம்} (caraṇālayam)}{சரணாலயம் (caraṇālayam)}]{\ta{சரணாலயம்} (caraṇālayam) — a sanctuary}

\label{lesson:TA-C437-caranalayam}



\begin{quote}
Before the new one: say the Tamil for coral, then the Tamil for wildlife.
\end{quote}

\subsection*{You'll want to know: \ta{சரணாலயம்}}
\textbf{\ta{சரணாலயம்}} — \emph{caraṇālayam} — \textquotedblleft{}a sanctuary\textquotedblright{}.

\begin{quote}
\ta{இந்தச்} \ta{சரணாலயத்தில்} \ta{பல} \ta{நாடுகளிலிருந்து} \ta{வந்த} \ta{பறவைகள்} \ta{இருக்கின்றன}.
\end{quote}

\begin{quote}
\emph{intac caraṇālayattil pala nāṭukaḷiliruntu vanta paṟavaikaḷ irukkiṉṟaṉa.}
\end{quote}

\begin{quote}
In this sanctuary there are birds that came from many countries.
\end{quote}

\ta{பறவைகள்} \ta{சரணாலயம்} is a bird sanctuary.

\begin{cousinweb}
\textbf{\ta{சரணாலயம்}} (\emph{caraṇālayam}) is built from two words borrowed from Sanskrit: \emph{śaraṇa}, \textquotedblleft{}shelter\textquotedblright{} (CDIAL 12326), + \emph{ālaya}, \textquotedblleft{}a house, a dwelling\textquotedblright{} (CDIAL 1366).
\end{cousinweb}

\begin{grammarlens}[title={What you've built}]
One more word for getting around.
\end{grammarlens}

\subsection*{Your turn}
\begin{itemize}
\raggedright
  \item \emph{Say it:} \emph{caraṇālayam}
  \item \emph{Say it:} \emph{caraṇālayam}, once more
  \item \emph{Say it:} \emph{intac caraṇālayattil pala nāṭukaḷiliruntu vanta paṟavaikaḷ irukkiṉṟaṉa.} — In this sanctuary there are birds that came from many countries.
  \item \emph{Recall:} say the Tamil for coral, then the Tamil for wildlife, then say \emph{caraṇālayam} again
\end{itemize}

\begin{tcolorbox}[breakable,colback=teal!4,colframe=teal!35!black,title={Before you move on}]
What does \ta{சரணாலயம்} mean? (\textquotedblleft{}A sanctuary\textquotedblright{}.) Say it once more.
\end{tcolorbox}
\section[\texorpdfstring{\ta{இடம்பெயர்} (iṭampeyar)}{இடம்பெயர் (iṭampeyar)}]{\ta{இடம்பெயர்} (iṭampeyar) — to migrate; to move away}

\label{lesson:TA-C437-itampeyar}



\begin{quote}
Before the new one: say the Tamil for wildlife, then the Tamil for a sanctuary.
\end{quote}

\subsection*{You'll want to know: \ta{இடம்பெயர்}}
\textbf{\ta{இடம்பெயர்}} — \emph{iṭampeyar} — \textquotedblleft{}to migrate; to move away\textquotedblright{}.

\begin{quote}
\ta{குளிர்} \ta{காலத்தில்} \ta{இந்தப்} \ta{பறவைகள்} \ta{தெற்கே} \ta{இடம்பெயர்கின்றன}.
\end{quote}

\begin{quote}
\emph{kuḷir kālattil intap paṟavaikaḷ teṟkē iṭampeyarkiṉṟaṉa.}
\end{quote}

\begin{quote}
In winter these birds migrate south.
\end{quote}

Past (I …): \ta{இடம்பெயர்ந்தேன்}, \emph{iṭampeyarntēṉ}.

\ta{இடம்}, a place, + \ta{பெயர்}, to move.

\begin{cousinweb}
\textbf{\ta{இடம்பெயர்}} (\emph{iṭampeyar}) is built from \ta{இடம்}, \textquotedblleft{}a place\textquotedblright{} (DEDR 434), and \ta{பெயர்}, \textquotedblleft{}to move, to depart\textquotedblright{} (DEDR 4409).
\end{cousinweb}

\begin{grammarlens}[title={What you've built}]
One more word for getting around.
\end{grammarlens}

\subsection*{Your turn}
\begin{itemize}
\raggedright
  \item \emph{Say it:} \emph{iṭampeyar}
  \item \emph{Say it:} \emph{iṭampeyar}, once more
  \item \emph{Say it:} \emph{kuḷir kālattil intap paṟavaikaḷ teṟkē iṭampeyarkiṉṟaṉa.} — In winter these birds migrate south.
  \item \emph{Recall:} say the Tamil for wildlife, then the Tamil for a sanctuary, then say \emph{iṭampeyar} again
\end{itemize}

\begin{tcolorbox}[breakable,colback=teal!4,colframe=teal!35!black,title={Before you move on}]
What does \ta{இடம்பெயர்} mean? (\textquotedblleft{}To migrate; to move away\textquotedblright{}.) Say it once more.
\end{tcolorbox}
\section[\texorpdfstring{\ta{வேட்டையாடு} (vēṭṭaiyāṭu)}{வேட்டையாடு (vēṭṭaiyāṭu)}]{\ta{வேட்டையாடு} (vēṭṭaiyāṭu) — to hunt}

\label{lesson:TA-C437-vettaiyatu}



\begin{quote}
Before the new one: say the Tamil for a sanctuary, then the Tamil for to migrate.
\end{quote}

\subsection*{You'll want to know: \ta{வேட்டையாடு}}
\textbf{\ta{வேட்டையாடு}} — \emph{vēṭṭaiyāṭu} — \textquotedblleft{}to hunt\textquotedblright{}.

\begin{quote}
\ta{புலி} \ta{இரவில்} \ta{வேட்டையாடுகிறது}.
\end{quote}

\begin{quote}
\emph{puli iravil vēṭṭaiyāṭukiṟatu.}
\end{quote}

\begin{quote}
The tiger hunts at night.
\end{quote}

Past (I …): \ta{வேட்டையாடினேன்}, \emph{vēṭṭaiyāṭiṉēṉ}.

\ta{வேட்டை}, hunting, + \ta{ஆடு}, a verb-making second part.

\begin{cousinweb}
\textbf{\ta{வேட்டையாடு}} (\emph{vēṭṭaiyāṭu}) is built from \ta{வேட்டை}, \textquotedblleft{}hunting, the chase\textquotedblright{} (DEDR 5527), and \ta{ஆடு}, \textquotedblleft{}to move, to play\textquotedblright{} (DEDR 347).
\end{cousinweb}

\begin{grammarlens}[title={What you've built}]
One more word for getting around.
\end{grammarlens}

\subsection*{Your turn}
\begin{itemize}
\raggedright
  \item \emph{Say it:} \emph{vēṭṭaiyāṭu}
  \item \emph{Say it:} \emph{vēṭṭaiyāṭu}, once more
  \item \emph{Say it:} \emph{puli iravil vēṭṭaiyāṭukiṟatu.} — The tiger hunts at night.
  \item \emph{Recall:} say the Tamil for a sanctuary, then the Tamil for to migrate, then say \emph{vēṭṭaiyāṭu} again
\end{itemize}

\begin{tcolorbox}[breakable,colback=teal!4,colframe=teal!35!black,title={Before you move on}]
What does \ta{வேட்டையாடு} mean? (\textquotedblleft{}To hunt\textquotedblright{}.) Say it once more.
\end{tcolorbox}
\section[\texorpdfstring{\ta{இனம்} (iṉam)}{இனம் (iṉam)}]{\ta{இனம்} (iṉam) — a species; a kind}

\label{lesson:TA-C437-inam}



\begin{quote}
Before the new one: say the Tamil for to migrate, then the Tamil for to hunt.
\end{quote}

\subsection*{You'll want to know: \ta{இனம்}}
\textbf{\ta{இனம்}} — \emph{iṉam} — \textquotedblleft{}a species; a kind\textquotedblright{}.

\begin{quote}
\ta{இந்தப்} \ta{பறவை} \ta{இனம்} \ta{இப்போது} \ta{மிக} \ta{அபூர்வம்}.
\end{quote}

\begin{quote}
\emph{intap paṟavai iṉam ippōtu mika apūrvam.}
\end{quote}

\begin{quote}
This species of bird is now very rare.
\end{quote}

Also a group of people or things of one kind.

\begin{cousinweb}
\textbf{\ta{இனம்}} (\emph{iṉam}) is native Dravidian (DEDR 531), \textquotedblleft{}a class, a group, a kind, a species\textquotedblright{}.
\end{cousinweb}

\begin{grammarlens}[title={What you've built}]
One more word for getting around.
\end{grammarlens}

\subsection*{Your turn}
\begin{itemize}
\raggedright
  \item \emph{Say it:} \emph{iṉam}
  \item \emph{Say it:} \emph{iṉam}, once more
  \item \emph{Say it:} \emph{intap paṟavai iṉam ippōtu mika apūrvam.} — This species of bird is now very rare.
  \item \emph{Recall:} say the Tamil for to migrate, then the Tamil for to hunt, then say \emph{iṉam} again
\end{itemize}

\begin{tcolorbox}[breakable,colback=teal!4,colframe=teal!35!black,title={Before you move on}]
What does \ta{இனம்} mean? (\textquotedblleft{}A species; a kind\textquotedblright{}.) Say it once more.
\end{tcolorbox}
